\begin{table}[htbp]
	\centering
	% 1. 标题设置
	\caption*{\zihao{3}\bfseries\heiti 社会调查小组成员评分表}
	
	% 2. 表格参数设置
	\renewcommand{\arraystretch}{0.6904} % 稍微调整行高，平衡内容密度
	\setlength{\tabcolsep}{4.6pt}       % 稍微收紧列间距
	
	% 3. 表格定义
	\begin{xltabular}{\textwidth}{|c|c|X|c|c|}
		\hline
		
		% --- 表头 ---
		\multirow{2}{*}{\textbf{序号}} & 
		\multirow{2}{*}{\textbf{姓名}} & 
		\multicolumn{1}{c|}{\multirow{2}{*}{\textbf{任务分工及答辩内容}}} & % 修改了表头名称使其更准确
		\multicolumn{2}{c|}{\textbf{评\quad 分}} \\
		\cline{4-5} 
		& & & \textbf{满分} & \textbf{得分} \\
		\hline
		
		% --- 第 1 位成员：侧重建模与分析 ---
		\multirow{2}{*}{1} & 
		\multirow{2}{*}{\parbox{4em}{\centering 杜浩月}} & 
		(1) 负责核心分析模型的构建与调试，运用统计软件进行数据挖掘与假设检验，确保分析结果的可靠性。 & 
		\makecell{表现 \\ 20} & 
		\multirow{2}{*}{} \\ 
		\cline{4-4} 
		& & 
		(2) 重点阐述模型选择的依据、技术路径及数据分析结果的统计学意义，回答关于数据处理细节的技术性提问。& 
		\makecell{答辩 \\ 30} & 
		\\ 
		\hline
		
		% --- 第 2 位成员：侧重写作与整合 ---
		\multirow{2}{*}{2} & % 修正序号为 2
		\multirow{2}{*}{\parbox{4em}{\centering 陈锦湘}} & 
		(1)  主导调查报告的整体撰写与统稿，整合各板块内容，梳理论证逻辑，确保报告符合学术规范。 & 
		\makecell{表现 \\ 20} & 
		\multirow{2}{*}{} \\ 
		\cline{4-4} 
		& & 
		(2) 汇报研究结论、对策建议及研究不足，侧重阐述研究发现的社会现实意义，回应关于报告整体结构的质询。 & 
		\makecell{答辩 \\ 30} & 
		\\ 
		\hline
		
		% --- 第 3 位成员：侧重设计与呈现 ---
		\multirow{2}{*}{3} & % 修正序号为 3
		\multirow{2}{*}{\parbox{4em}{\centering 刘志琦}} & 
		(1)  负责调查问卷的结构化设计与题项编制，参与预调查及信效度初步检验。配合完成最终报告的图表绘制与排版美化。 & 
		\makecell{表现 \\ 20} & 
		\multirow{2}{*}{} \\ 
		\cline{4-4} 
		& & 
		(2)  介绍问卷的设计思路、核心变量的操作化定义过程，以及如何通过设计保证调查的科学性。 & 
		\makecell{答辩 \\ 30} & 
		\\ 
		\hline
		
		% --- 第 4 位成员：侧重执行与数据清洗 ---
		\multirow{2}{*}{4} & % 修正序号为 4
		\multirow{2}{*}{\parbox{4em}{\centering 王嘉伟}} & 
		(1)  执行问卷的分发策略，监控样本回收进度与质量。负责原始数据的清洗、编码、录入及基础描述性统计分析。 & 
		\makecell{表现 \\ 20} & 
		\multirow{2}{*}{} \\ 
		\cline{4-4} 
		& & 
		(2) 说明样本获取渠道的有效性、样本结构的代表性以及数据清洗所采用的标准和流程。 & 
		\makecell{答辩 \\ 30} & 
		\\ 
		\hline
		
		% --- 第 5 位成员：侧重统筹与展示 ---
		\multirow{2}{*}{5} & % 修正序号为 5
		\multirow{2}{*}{\parbox{4em}{\centering 马士豪}} & 
		(1)  统筹项目整体进度与分工协作，管理研究文献与过程资料。负责答辩 PPT 的逻辑框架搭建与核心内容提炼。 & 
		\makecell{表现 \\ 20} & 
		\multirow{2}{*}{} \\ 
		\cline{4-4} 
		& & 
		(2)  负责答辩的开场背景介绍与总结陈词，宏观把控答辩节奏，协调组员回答跨模块的综合性问题。& 
		\makecell{答辩 \\ 30} & 
		\\ 
		\hline
		
	\end{xltabular}
\end{table}